

There is abundant astrophysical evidence for the existence of dark matter~\cite{Aghanim2020, SofueRubin, harvey2015nongravitational, ARBEY2021103865}, a nonrelativistic and nonbaryonic matter component of the Universe that has so far eluded direct detection through interaction with ordinary matter~\cite{schumann2019direct}. Weakly interacting massive particles (WIMPs), which obtain their relic abundance by thermal freeze-out through weak interactions~\cite{Lee1977ua}, are postulated in a wide variety of viable extensions to the standard model of particle physics~\cite{RevModPhys.90.045002, Billard2022,akerib2022snowmass2021}. They are a leading candidate to explain dark matter, despite strong constraints from many searches completed and ongoing at colliders~\cite{aaboud2018search, sirunyan2018search, aaboud2018searchz, sirunyan2018searchz, universe4110131}, with telescopes~\cite{abe2020indirect,choi2015search, abbasi2022search, albert2017results,  cuoco2017novel, cui2017possible, Albert_2017}, and in underground laboratories~\cite{abdelhameed2019first, agnes2018low,amole2019dark,collaboration2018dark, PhysRevLett.127.261802,deapcollaborationSearchDarkMatter2019, akerib2017results, Agnese:2015ywx}. 
This Letter reports the first search for dark matter from the LUX-ZEPLIN (LZ) experiment, with the largest target mass of any WIMP detection experiment to date.

The LZ experiment~\cite{LZExperiment, lztdr} is located \SI{4850}{ft} underground in the Davis Cavern at the Sanford Underground Research Facility (SURF) in Lead, South Dakota, USA, shielded by an overburden of \SI{4300}{m} water-equivalent~\cite{Heise:2021eym}. It is a low-background, multidetector experiment centered on a dual-phase time projection chamber (TPC) mounted in a double-walled titanium cryostat~\cite{AKERIB20171} filled with \SI{10}{t} of liquid xenon (LXe). The TPC is a vertical cylinder approximately \SI{1.5}{m} in diameter and height, lined with reflective polytetrafluoroethylene, and instrumented with 494 3{\nobreakdash-}inch photomultiplier tubes (PMTs) in two arrays at top and bottom. Energy depositions above approximately \SI{1}{keV} in the \SI{7}{t} active xenon region produce two observable signals: vacuum ultraviolet (VUV) scintillation photons (S1) and ionization electrons that drift under a uniform electric field to the liquid surface, where they are extracted and produce secondary scintillation in the xenon gas (S2). The ratio of S2 to S1 differentiates interactions with a xenon nucleus (producing a nuclear recoil, or NR) from interactions with the atomic electron cloud (producing an electron recoil, or ER). 
 
The TPC is surrounded by two detectors, which provide veto signals to reject internal and external backgrounds. A LXe ``skin'' detector between the TPC field cage and the cryostat wall is instrumented with 93 1{\nobreakdash-}inch and 38 2{\nobreakdash-}inch PMTs.  The outer detector (OD) is a near-hermetic system of acrylic tanks containing \SI{17}{t} of gadolinium-loaded (\SI{0.1}{\percent} by mass) liquid scintillator~\cite{Haselschwardt:2018vmp} to detect neutrons. The entire LZ detector system is in a tank filled with \SI{238}{t} of ultrapure water to shield from the ambient radioactive background, and 120 8-inch PMTs are submersed in the water to record OD and water Cherenkov signals.
